\documentclass{article}
\usepackage[T1]{fontenc}
\usepackage{lmodern}
\usepackage{graphicx}
\usepackage{amsmath,amssymb}
\usepackage[a4paper,margin=2cm]{geometry}
\usepackage[absolute,overlay]{textpos}
\setlength{\TPHorizModule}{1mm}
\setlength{\TPVertModule}{1mm}
\usepackage[indonesian]{babel}
\usepackage{hyperref}
\setlength{\emergencystretch}{2em}
% unit: Halaman 11

\begin{document}

% === Halaman 11 (llm) ===

\begin{itemize}
    \item[$\bullet$] \textbf{Contoh 2:}
\end{itemize}

\begin{tabular}{ll}
    Plainteks: & \texttt{nantimalamsayatunggukamudidepanwarungkopi} \\
    Kunci: & \texttt{gtrskncvbrwpoatqljfmxr pjsrzfoltfbtmaedpvy} \\
    Cipherteks: & \texttt{TTELSZCGBDOPMAMKYPLGHTDJMAUDDLSDTSGNKNDKG}
\end{tabular}

\end{document}
